% \makeatletter
% \def\input@path{{eg-template}}
% \makeatother

\documentclass[twocolumn,a4paper]{article}
% \documentclass{egpubl}

% \usepackage{EGUKCGVC2023}
% \WsSubmission
\usepackage[T1]{fontenc}
\usepackage[utf8]{inputenc}
% \usepackage{dfadobe}

% \biberVersion
% \BibtexOrBiblatex
\usepackage[backend=biber,citestyle=alphabetic,backref=true]{biblatex}
\addbibresource{references.bib}
% \electronicVersion
% \PrintedOrElectronic

% \ifpdf \usepackage[pdftex]{graphicx} \pdfcompresslevel=9
% \else \usepackage[dvips]{graphicx} \fi
% \usepackage{egweblnk}
\usepackage{hyperref}

%%%
\usepackage{adjustbox}
\usepackage{amsmath}
\usepackage{amssymb}
% \usepackage[backend=biber,bibstyle=EG,citestyle=alphabetic]{biblatex}
\usepackage{float}
\usepackage{lipsum}
\usepackage{subcaption}
\usepackage{subfiles}
\usepackage{tabularx}
\usepackage[normalem]{ulem}
\usepackage{xcolor}
\usepackage{xurl}

\graphicspath{{./figures}}
\newcommand{\orbit}[1]{%
    \langle#1\rangle
}

\newcommand{\PN}[1]{%
    \{#1\}
}

\newcommand{\commentout}[1]{}

\newcommand{\maxime}[1]{%
    \textbf{\color{blue} Maxime: #1}
}

\newcommand{\agnes}[1]{%
    \textbf{\color{red} Agnès: #1}
}

\newcommand{\david}[1]{%
    \textbf{\color{violet} David: #1}
}

\newcommand{\stephane}[1]{%
    \textbf{\color{green} [Stephane]: #1}
}
\newcommand{\xavier}[1]{%
    \textbf{\color{magenta} [Xavier]: #1}
}
\newcommand{\xavierSupp}[1]{%
    \textbf{\color{magenta} \sout{#1}}
}
\newcommand{\xavierAdd}[1]{%
    \textbf{\color{magenta} #1}
}

\newcommand{\hakim}[1]{%
    \textbf{\color{teal} [Hakim]: #1}
}

\definecolor{darkgreen}{rgb}{0,.4,0}

\newcommand{\highlight}[1]{%
    \textbf{\color{orange}#1}
}

\newenvironment{proofidea}{\noindent\textbf{Idea of proof:}}{$\square$ \newline}
\newtheorem{definition}{Definition}[subsection]
\newtheorem{proposition}{Proposition}[subsection]



\title{Model reevaluation based on graph transformation rules}
\author{Anonymous}
\date{\today}

\begin{document}

\maketitle


\begin{abstract}
   % The ABSTRACT is to be in fully-justified italicized text,
   % between two horizontal lines,
   % in one-column format,
   % below the author and affiliation information.
   % Use the word ``Abstract'' as the title, in 9-point Times, boldface type,
   % left-aligned to the text, initially capitalized.
   % The abstract is to be in 9-point, single-spaced type.
   % The abstract may be up to 3 inches (7.62 cm) long. \\
   % Leave one blank line after the abstract,
   % then add the subject categories according to the ACM Classification Index
In this paper, we widen the naming problem studies to the rule-based graph transformation modeling systems. 
We propose a persistent naming method taking advantage of the generalized maps' and graph transformation rules’ formalization of operations.
It enables a unique and homogeneous characterisation of entities in all dimensions.
Most existing methods require tracking numerous topological entities and consider the persistent naming problem only from the parameters' modifications of a parametric specification standpoint. 
%In 
With 
our solution, not only the naming problem is tackled within the usual framework of parameters edition, but we also take the specification edition into account %(operations addition, deletion, move).
(addition, deletion and displacement of operations).
Moreover, our solution makes use of directed acyclic graphs %(DAGs) 
to represent the histories of 
topological entities%' histories 
and to track only the entities used in the parametric specification and the ones they originate from. 
Representing 
the
history of a topological entity 
%both in an evaluation DAG and after editing the parametric specification in a reevaluation DAG 
allows for an efficient persistent naming mechanism that takes into account the impact
%of both origins and traces
of the modifications during reevaluation.
%-------------------------------------------------------------------------
%  ACM CCS 1998
%  (see http://www.acm.org/about/class/1998)
% \begin{classification} % according to http:http://www.acm.org/about/class/1998
% \CCScat{Computer Graphics}{I.3.3}{Picture/Image Generation}{Line and curve generation}
% \end{classification}
%-------------------------------------------------------------------------
%  ACM CCS 2012
%  (see http://www.acm.org/about/class/class/2012)
%The tool at \url{http://dl.acm.org/ccs.cfm} can be used to generate
% CCS codes.
%Example:
% \begin{CCSXML}
% <ccs2012>
%    <concept>
%        <concept_id>10010147.10010371.10010396</concept_id>
%        <concept_desc>Computing methodologies~Shape modeling</concept_desc>
%        <concept_significance>500</concept_significance>
%        </concept>
%  </ccs2012>
% \end{CCSXML}

% \ccsdesc[500]{Computing methodologies~Shape modeling}
% \printccsdesc

\textbf{Keywords}: Topology-based modeling; Graph transformation rules; Persistent Naming; Reevaluation; Generalized maps;
\end{abstract}

\subfile{sections/introduction.tex}
\subfile{sections/main-concepts.tex}
\subfile{sections/parametric-specifications.tex}
\subfile{sections/evaluation.tex}
\subfile{sections/reevaluation.tex}
%\subfile{sections/implementation.tex}
\subfile{sections/conclusions-persp.tex}

\printbibliography


\end{document}
